\section{Introdução}

Doenças cardiovasculares (DCV) são a principal causa de morte no mundo. A organização mundial da saúde (OMS) estima que em 2030 quase 23,6 milhões de pessoas morrerão de doenças cardiovasculares \cite{world2011global}. Na população pediátrica o risco de doenças cardiovasculares estão mais associados a obesidade e sedentarismo \cite{dias_risco_2023}, mas existem outras condições que podem levar a doenças cardiovasculares, como a cardiopatia congênita, que é uma das principais causas de morte em crianças \cite{born20098}.

Os métodos de diagnóstico atuais para doenças cardiovasculares incluem exames de imagem como eletrocardiograma (ECG), ecocardiograma, exame radiológico do torax, ressonância magnética, tomografia computadorizada, entre outros. Alem disso dispositivos tecnológicos como, relogios e aneis inteligentes estão sendo amplamente adotados para monitorar estilo de vida e fatores de risco cardiovascular \cite{akinosun2021digital}. Apesar da diversidade de exames disponíveis, os custos proibitivos, bem como a falta de concientização pública sobre os fatores de risco cardiovascular e falta de exames preventivos regulares são os pricipais fatores que contibuem para o diagnóstico tardio de doenças cardiovasculares \cite{magalhaes2024desafios}.

Dado a severidade das doenças cardiovasculares, a detecção precoce é crucial para o tratamento e prevenção de complicações. A inteligência artificial (IA) tem sido amplamente utilizada em diversas áreas da saúde, incluindo diagnóstico de doenças cardiovasculares \cite{subramani_cardiovascular_2023,dahia_implementing_2023, abdulsalam2023explainable,linardos2022federated}. A adoção da IA advém principalmente da capacidade de processar grandes volumes de dados e identificar padrões complexos que podem não ser facilmente identificados por humanos \cite{van2022critical}.

Neste contexto, esta pesquisa avalia diferentes técnicas de aprendizado de máquina e redes neurais artificiais para a detecção de doenças cardiovasculares em crianças. O objetivo é identificar a técnica mais adequada para a detecção de doenças cardiovasculares em crianças, considerando a precisão, sensibilidade e especificidade do modelo.

Esta pesquisa está estruturada da seguinte forma: a seção 2 apresenta a revisão da literatura sobre o uso de inteligência artificial para detecção de doenças cardiovasculares em crianças. As seções 3 e 4 apresentam a metodologia utilizada para a coleta de dados, pré-processamento, treinamento e avaliação dos modelos. A seção 5 apresenta os resultados obtidos, a seção 6 apresenta a estratégia final para a validação em ambiente de produção e a seção 7 apresenta as considerações finais.